\section{Limitations}
\label{sec:limitations}
\noindent We highlight key limitations and threats to validity of this study.

\begin{enumerate}
    \item Our graphs are centered on T2DM ($\mathbb{G}_1$), AD ($\mathbb{G}_2$) and AD+T2DM ($\mathbb{G}_3$) space. Findings might not be related to other conditions, specialities, non-PubMed corpora, and multilingual environments

    \item Although co-reference and canonicalization have been improved, the errors in entity connecting, synonym merging, and relation extracting may be carried over to the retrieval to produce plausible yet off-target evidence that diminishes accuracy. We used rule-based filtering to reduce false positives, but the graph is not 100\% perfect.


    \item We did not actively search retriever depth, hybrid lexical neural retrieval or learning-to-rank rerankers. A more powerful ranking stack that has the potential to minimize distractors, particularly on union graphs.

    \item Accuracy, Macro P/R/F1, and Micro F1 fail to reflect on calibration, factual grounding to primary sources and clinical harm potential. The fact-checking based on human judgment and reference was out of the question.


    \item The findings on smaller vs. larger models might not be generalizable to other architectures, tokenizer selection, alignment processes or domain-trained checkpoints.

    \item Our setup omits latency, cost, privacy, and PHI-handling constraints crucial in clinical workflows in this pilot study. Deployment-time retrieval drift, updates to literature, and governance requirements are not modeled.

    \item Since we had to run on a few resources, we did not carry out large multi-seed reruns, ablations of KG preprocessing steps, or confidence-interval reporting over all settings; here, some of the effects may be due to stochasticity.
\end{enumerate}
